% \chapter{Messdaten}
\setcounter{chapter}{2}
\addcontentsline{toc}{chapter}{\thechapter \, Messdaten} % Damit es im Inhaltsverzeichnis erscheint
\setcounter{table}{0}  % Setze den Tabellenzähler zurück
\label{ch:Protokoll}

\thispagestyle{empty}

\includepdf[
  pages=1-8,               
  pagecommand={\thispagestyle{empty}} 
]{Versuche/\versuchsnummer/Auswertung/Messprotokoll.pdf}



\addcontentsline{lot}{table}{{M.1.1} Halbwertszeiten Spannung}
\addcontentsline{lot}{table}{{M.1.2} Halbwertszeiten Strom}
\addcontentsline{lot}{table}{{M.3.1} Grenzfrequenzen}
\addcontentsline{lot}{table}{{M.3.2} Phasenverschiebung Tiefpass}
\addcontentsline{lot}{table}{{M.3.3} Phasenverschiebung Hochpass}
\addcontentsline{lot}{table}{{M.4.1} Grenzfrequenzen RLC}
\addcontentsline{lot}{table}{{M.5.1} Resonantfrequenzen}
\addcontentsline{lot}{table}{{M.6.1} Spannungsamplitude und Periodendauer}
\addcontentsline{lot}{table}{{M.8.1 - M.8.6} Amplituden}